\section{Discussion}
This study shows that one central EEG channel can detect deep sleep and describe it in physiological terms. Four findings stand out. First, a compact detector matched DeepSleepNet, which has 264 times more parameters, and outperformed three other re-implemented architectures under the same protocol, mainly thanks to temporal context \cite{phan2022sleeptransformer,phan2023lseqsleepnet}. Second, it ranked N3 epochs almost as well in five independent cohorts as in cross-validation, and where experts agreed least, its agreement with them was within their own range. Third, the transparent scores followed known slow-wave physiology: they decreased across the night \cite{borbely2016twoprocess,riedner2007slowwaves} and with age, and rose with the share of experts who scored N3. Fourth, the reliability score identified the epochs that need review. % rq1_sota_paired_tests, rq4_component_effects, rq1_interrater_benchmark

\subsection{Implications for long-term monitoring of deep sleep}
Large studies of slow-wave sleep and brain health \cite{xie2013clearance,baril2021brainaging,himali2023dementia} need measures that are cheap, repeatable and understandable. Single-channel home EEG can match polysomnography for slow-wave activity \cite{lucey2016singlechannel}, but wearables rarely measure slow-wave sleep accurately \cite{green2022devices}. Our method needs one central channel and 96,226 parameters, reports only physiologically named numbers with an exact additive breakdown, and gives every epoch a reliability score. Excluding the least reliable fifth raised F1 on the remaining epochs to 0.949, but reliability was weaker in DOD (area 0.751 and 0.689), where the detector is least accurate, so any monitoring use would need a check of this kind before its output is trusted. % rq6_risk_coverage

\subsection{Interpreting the Dreem results}
\label{sec:dod}
The DOD results are the main limitation of the detector. The evidence suggests that the problem is where the reference places the N2/N3 boundary, not a failure to see slow waves. The ranking of epochs and the order of nights transferred to DOD (Fig.~\ref{fig:ba}b). False positives were mostly consensus N2 epochs with less automatically detected slow-wave coverage than true positives (median 4.6\% against 18.5\% in DOD-O). This fits DOD scorers placing the boundary higher, although automated coverage is generally lower than visual coverage. The same pattern appeared without class weighting (DOD-O precision 0.415), and YASA showed the same precision drop in DOD; its more cautious operating point gave a higher DOD F1 but a lower F1 elsewhere. A threshold chosen for each cohort on test labels would raise the DOD F1 score, but it would make the external evaluation invalid. % rq1_detection_metrics; sensitivity_rq1_false_positive_structure; rq4_variant_external_metrics; rq7_profile_agreement icc31

Calibration showed a related effect (Section~IV-A). Positive-class weighting can move the logits by a constant amount that a single temperature cannot remove \cite{guo2017calibration}. Without class weighting, the cross-validated error fell from 0.060 to 0.006, and Platt scaling of the primary model reached 0.003; in DOD both corrections reduced it to 0.11--0.15, and in the other external cohorts it stayed at or below 0.044. This is consistent with the logit-shift explanation, the remaining DOD error fits the boundary shift, and Platt scaling is a simple correction in practice. % posthoc_calibration.csv

\subsection{What the factor scores measure}
The factor scores describe slow-wave strength in one EEG derivation and are meant as research measures. They are not validated measures of clinical, subjective or restorative sleep quality, because none of our datasets contains such a reference. Their agreement with the graded expert vote is convergent validity: the scores measure, on a continuous scale, the slow activity that makes experts vote N3, and so turn expert disagreement into a graded quantity. Because morphology uses the amplitude-and-coverage definition that experts apply, part of this agreement follows from the definition, so it is a consistency check rather than independent validation. The approach is in line with the hypnodensity view of stage ambiguity \cite{bakker2023ambiguity,huijben2023hypnodensity} and with concerns about fixed N3 criteria \cite{davidson2025n3}. The exploratory post hoc comparison suggests that the two outputs play complementary roles. The detector logit, trained on expert labels, matched expert judgement best and is the better number for ranking epochs as experts would. Slow-wave morphology followed the homeostatic decline, which does not depend on expert scoring, more strongly than the logit, also on epochs not selected by the detector, and delta intensity did so in most comparisons. The scores also decreased with age, as slow-wave amplitude and other EEG features do \cite{rosinvil2021age,mai2025aisleep}, even after adjustment for diagnosis, and, unlike the logit, they are named, additive and comparable between nights. The combined rating matched the logit on the decline. The test--retest data from eight participants are consistent with stability but too small to confirm it. % rq2_factor_validity graded; posthoc_icc_ci

Because the factors are strongly correlated, the rating hardly depends on the prior weights, and its exact additive breakdown shows which factor caused a given value. % rq3_overall_rating_stability.json

\subsection{The slow-wave grade}
The grade came close to its criterion (0.743 against 0.75) and was cautious: almost no stage 3 epoch was graded High, but only one in five stage 4 epochs was. A likely reason is that automated coverage from a filtered channel is lower than a scorer's visual estimate. The thresholds come from both scoring systems and the reference is R\&K-scored CAP, so this result does not establish agreement with AASM grading. A lower High threshold would have to be learned from the test reference, so we report the grade as not yet validated; the continuous rating (0.776) and the detector logit (0.860) separated stage 4 from stage 3 better. % rq2_rk_grade_validation; sensitivity_rq2b

\subsection{Reliability and explanations}
The reliability score separates likely errors well in the training cohorts, ISRUC and CAP, and it identifies epochs on which human experts disagree. This matches what ensemble and Monte Carlo uncertainty achieve for five-class staging \cite{dei2025sleepyland,vainikka2026uncertainty}. The quality-control flags did not predict errors because the detector makes its mistakes on clean N2/N3 boundary epochs, not on artefacts. In DOD, error detection was weaker, as expected for a systematic boundary shift, although the out-of-distribution and quality components helped most there. Both explanation methods ranked the delta band first in all analysed epochs, including errors. Because N3 is defined by delta waves, this is consistent with the physiology that defines N3, but it does not by itself show that the network learned slow waves. % rq6_reliability_components; sensitivity_rq6; rq5x

\subsection{Limitations}
\label{sec:limitations}
First, public data contain no expert ratings of slow-wave strength, so the factor scores and grade are research descriptors tested only against indirect references; a blinded annotation of within-N3 slow-wave strength by several raters is the natural next step, and the open code makes such a study straightforward. Second, in ISRUC-III the amplitude-dependent factors were close to the floor even for expert N3 epochs (median delta-intensity percentile 4.6 against 44.2 in ISRUC-I), while the scale-free spectral factor was normal (47.1). This points to an amplitude-scale difference in the source files of this cohort; the amplitude-dependent YASA also found only 43\% of expert N3 epochs there. Detection, which uses recording-normalised input, and the within-participant analyses are not affected, but absolute scores and grades in ISRUC-III cannot be interpreted. Third, individual age and sex are available only for CAP. There, the age association held after adjustment for diagnosis, but because the scores depend on age and sex \cite{davidson2025n3,rosinvil2021age}, absolute values can only be compared within a similar population and recording set-up, and age-specific references are needed. Fourth, all cohorts used a central-mastoid derivation; frontal and wearable montages need their own percentile references and evaluation. Fifth, the test--retest analysis is based on only eight participants. Sixth, the baselines were re-implemented with a shared recipe (Table~SI) and may perform below their best published configurations. Seventh, the DOD labels needed an alignment step that uses labels, and a label-free alignment has not yet been tested. Eighth, DOD-O is an obstructive sleep apnoea population whose arousals and architecture differ from the training cohorts. Finally, the post hoc analyses do not affect any verdict; Table~SVI lists where their conclusions differed. % ledger_extra 4.6, 44.2, 47.1
